\chapter*{Vorwort}
\addcontentsline{toc}{chapter}{Vorwort}
\markboth{Vorwort}{Vorwort}

\begin{tcolorbox}[colback=infobg, colframe=primary, arc=6pt, boxrule=0.8pt,
                  left=10pt, right=10pt, top=10pt, bottom=10pt]
  \faQuoteLeft\quad\textit{Gutes Training entsteht nicht im Moment des Kurses –
  es entsteht in der Vorbereitung, im Austausch und in der gemeinsamen
  Weiterentwicklung des Teams.}
\end{tcolorbox}

\bigskip

Dieses Buch ist ein internes Arbeitsmittel des \textbf{Dog College Franken}
und richtet sich ausschließlich an unser Trainerteam. Es sammelt bewährte
Übungen aus unserer täglichen Kursarbeit und soll als gemeinsame Grundlage
dienen – für erfahrene Trainer ebenso wie für neue Kollegen, die neu zu
uns stoßen.

\bigskip

\textbf{Wozu dieses Buch?}

\begin{itemize}[leftmargin=2em, itemsep=6pt]
  \item[\faLightbulb] \textbf{Inspiration} – Wer schon lange unterrichtet,
    kennt das Gefühl: Man greift immer wieder auf die gleichen Übungen zurück.
    Dieses Buch soll neue Impulse geben und das eigene Repertoire erweitern.
  \item[\faUserTie] \textbf{Einarbeitung} – Neue Trainer finden hier einen
    strukturierten Einstieg in unsere Kursinhalte und unsere Art zu arbeiten.
  \item[\faUsers] \textbf{Einheitlichkeit} – Damit Teilnehmer, die zwischen
    Kursen oder Trainern wechseln, ein konsistentes Training erleben.
\end{itemize}

\bigskip

\textbf{Aufbau des Buches}

Das Buch gliedert sich in zwei Hauptteile:

\begin{itemize}[leftmargin=2em, itemsep=6pt]
  \item[\faHome] \textbf{Heimübungen} – Übungen, die Hundehalter eigenständig
    zu Hause durchführen können. Ideal als Hausaufgaben zwischen den Kursstunden.
  \item[\faUsers] \textbf{Gruppenübungen} – Übungen, die die Gruppensituation
    gezielt als Trainingsreiz nutzen und im Kurs am wirkungsvollsten sind.
\end{itemize}

\smallskip

Innerhalb jedes Teils sind die Übungen nach Altersgruppen geordnet:
\textbf{Welpen \& Junghunde}, \textbf{Pubertät} und \textbf{Erwachsene Hunde}.
Jede Übungsseite enthält neben der Beschreibung auch einen Abschnitt
\textit{Variationen} sowie \textit{Trainer-Hinweise} – dort sind besondere
Beobachtungspunkte und häufige Fehler festgehalten.

\bigskip

Dieses Buch lebt vom Mitmachen: Wer eine Übung weiterentwickelt, eine neue
Idee hat oder einen hilfreichen Trainer-Tipp entdeckt – her damit. Es soll
wachsen und sich mit unserem Team weiterentwickeln.

\bigskip

\begin{flushright}
  \textit{Viel Freude beim Stöbern und Ausprobieren!}\\[6pt]
  \textit{Euer Dog College Franken Team}\\[4pt]
  \textcolor{primary}{\faPaw~\faPaw~\faPaw}
\end{flushright}

\clearpage
